\documentclass[tikz,border=8pt]{standalone}
\usepackage{tikz}
\usepackage{amsmath}
\usepackage{amssymb}
\usepackage{ctex}
\usetikzlibrary{calc, positioning, arrows.meta, decorations.pathreplacing, fit, backgrounds}

\begin{document}
\begin{tikzpicture}[>=Stealth,
  dombox/.style={draw=gray!50, fill=gray!5, rounded corners=8pt,
                 minimum width=3.6cm, minimum height=5.0cm},
  subbox/.style={draw=#1!60, fill=#1!10, rounded corners=5pt,
                 minimum width=3.0cm, minimum height=1.3cm,
                 inner sep=6pt, align=center, font=\small},
]

%% ===== 左侧大框：R^n（定义域）=====
\node[dombox] (Rn) at (0,0) {};
\node[font=\normalsize\bfseries, gray!60] at (0, 2.2) {$\mathbb{R}^n$};

% ker T（核）：映到 0 的子集
\node[subbox=orange] (ker) at (0, 0.65)
  {\textbf{核} $\ker T$\\所有映到 $\mathbf{0}$ 的向量\\$\dim = n - r$};

% 其余向量（补空间示意）
\node[subbox=blue] (comp) at (0, -1.1)
  {\textbf{补空间}\\（行空间方向）\\$\dim = r$};

%% ===== 中间：映射箭头 =====
% 核到 0
\draw[->, orange!70, thick, bend left=15]
  (ker.east) to node[above, font=\scriptsize, orange!80!black, sloped] {$T(\mathbf{x})=\mathbf{0}$}
  (4.5, 0.65);

% 补空间到像（双射）
\draw[->, blue!70, thick, bend right=15]
  (comp.east) to node[below, font=\scriptsize, blue!80!black, sloped] {双射}
  (4.5, -1.1);

% T 标签
\node[font=\Large\bfseries, gray!70] at (3.5, 0) {$T$};
\draw[->, very thick, gray!50] (2.0, 0) -- (4.8, 0);

%% ===== 右侧大框：R^m（陪域）=====
\node[dombox] (Rm) at (7.0, 0) {};
\node[font=\normalsize\bfseries, gray!60] at (7.0, 2.2) {$\mathbb{R}^m$};

% 零向量（核映到此）
\node[draw=orange!60, fill=orange!10, circle, minimum size=0.55cm,
      font=\scriptsize, align=center] (zero) at (7.0, 0.65) {$\mathbf{0}$};

% im T（像）
\node[subbox=teal] (im) at (7.0, -1.1)
  {\textbf{像} $\operatorname{im} T$\\$T$ 的所有输出\\$\dim = r$};

% 陪域其余部分标注
\node[font=\scriptsize, gray!50, align=center] at (7.0, 2.0) {};

%% ===== 维数定理标注框（外侧）=====
\node[font=\small, fill=white, draw=gray!50, thin, rounded corners=4pt,
      inner sep=8pt, align=center]
  at (3.5, -3.6) {
  \textbf{维数定理（秩-零化度定理）：}\\[4pt]
  $\dim(\ker T) + \dim(\operatorname{im} T) = n$\\[2pt]
  {\footnotesize 即 $(n - r) + r = n$，\quad $r = \operatorname{rank}(T)$}
};

%% ===== 整体标题 =====
\node[font=\large\bfseries] at (3.5, 3.5) {线性映射的核与像};

\end{tikzpicture}
\end{document}
